% GENERATED FILE. Edit canonical lessons, then run npm run generate:books.
% canonical-source-hash: fnv1a64:baf68a1250c88e70
% canonical-lessons: BN-C207-dam, BN-C207-dokandar, BN-C207-khodder, BN-C207-bbsa, BN-C207-kompani

\chapter{Price, Shopkeeper, Customer, Business, Company}
\label{ch:bn-price-shopkeeper-customer-business-company}
\hlchaptermodality{207}

\begin{chapteropening}
\textbf{By the end of this chapter:} \emph{I can say a price, a shopkeeper, a customer, business and a company.}

Complete the last lesson of chapter 207: I can say a price, a shopkeeper, a customer, business and a company.
\end{chapteropening}

\section[\texorpdfstring{\bn{দাম} (dām)}{দাম (dām)}]{\bn{দাম} (dām) — a price}

\label{lesson:BN-C207-dam}



\begin{quote}
Before the new one: say the Bengali for a tray, then the Bengali for a strainer.
\end{quote}

\subsection*{You'll want to know: \bn{দাম}}
\textbf{\bn{দাম}} — \emph{dām} — \textquotedblleft{}a price\textquotedblright{}.

\begin{cousinweb}
\textbf{\bn{দাম}} (\emph{dām}) is a tadbhava, a word inherited from Sanskrit and reshaped by sound change. It comes from \emph{dramma}, \textquotedblleft{}a coin\textquotedblright{} (CDIAL 6622). CDIAL lists Bengali \emph{dām}, \textquotedblleft{}price\textquotedblright{}, and takes \emph{dramma} from Greek \emph{drachmē} through Middle Persian.
\end{cousinweb}

\begin{grammarlens}[title={What you've built}]
One more word for this chapter.
\end{grammarlens}

\subsection*{Your turn}
\begin{itemize}
\raggedright
  \item \emph{Say it:} \emph{dām}
  \item \emph{Say it:} \emph{dām}, once more
  \item \emph{Say it:} say \emph{dām}
  \item \emph{Recall:} say the Bengali for a tray, then the Bengali for a strainer, then say \emph{dām} again
\end{itemize}

\begin{tcolorbox}[breakable,colback=teal!4,colframe=teal!35!black,title={Before you move on}]
What does \bn{দাম} mean? (\textquotedblleft{}A price\textquotedblright{}.) Say it once more.
\end{tcolorbox}
\section[\texorpdfstring{\bn{দোকানদার} (dokāndār)}{দোকানদার (dokāndār)}]{\bn{দোকানদার} (dokāndār) — a shopkeeper}

\label{lesson:BN-C207-dokandar}



\begin{quote}
Before the new one: say the Bengali for a strainer, then the Bengali for a price.
\end{quote}

\subsection*{You'll want to know: \bn{দোকানদার}}
\textbf{\bn{দোকানদার}} — \emph{dokāndār} — \textquotedblleft{}a shopkeeper\textquotedblright{}.

\begin{cousinweb}
\textbf{\bn{দোকানদার}} (\emph{dokāndār}) is \emph{dokān}, \textquotedblleft{}shop\textquotedblright{}, from Arabic \emph{dukkān}, which reached Bengali through Persian (Lane), + the Persian ending \emph{-dār}, \textquotedblleft{}holder\textquotedblright{}.
\end{cousinweb}

\begin{grammarlens}[title={What you've built}]
One more word for this chapter.
\end{grammarlens}

\subsection*{Your turn}
\begin{itemize}
\raggedright
  \item \emph{Say it:} \emph{dokāndār}
  \item \emph{Say it:} \emph{dokāndār}, once more
  \item \emph{Say it:} say \emph{dokāndār}
  \item \emph{Recall:} say the Bengali for a strainer, then the Bengali for a price, then say \emph{dokāndār} again
\end{itemize}

\begin{tcolorbox}[breakable,colback=teal!4,colframe=teal!35!black,title={Before you move on}]
What does \bn{দোকানদার} mean? (\textquotedblleft{}A shopkeeper\textquotedblright{}.) Say it once more.
\end{tcolorbox}
\section[\texorpdfstring{\bn{খদ্দের} (khôdder)}{খদ্দের (khôdder)}]{\bn{খদ্দের} (khôdder) — a customer}

\label{lesson:BN-C207-khodder}



\begin{quote}
Before the new one: say the Bengali for a price, then the Bengali for a shopkeeper.
\end{quote}

\subsection*{You'll want to know: \bn{খদ্দের}}
\textbf{\bn{খদ্দের}} — \emph{khôdder} — \textquotedblleft{}a customer\textquotedblright{}.

\begin{cousinweb}
\textbf{\bn{খদ্দের}} (\emph{khôdder}) is usually traced to Persian \emph{kharīdār}, \textquotedblleft{}buyer\textquotedblright{}, worn down in Bengali.
\end{cousinweb}

\begin{grammarlens}[title={What you've built}]
One more word for this chapter.
\end{grammarlens}

\subsection*{Your turn}
\begin{itemize}
\raggedright
  \item \emph{Say it:} \emph{khôdder}
  \item \emph{Say it:} \emph{khôdder}, once more
  \item \emph{Say it:} say \emph{khôdder}
  \item \emph{Recall:} say the Bengali for a price, then the Bengali for a shopkeeper, then say \emph{khôdder} again
\end{itemize}

\begin{tcolorbox}[breakable,colback=teal!4,colframe=teal!35!black,title={Before you move on}]
What does \bn{খদ্দের} mean? (\textquotedblleft{}A customer\textquotedblright{}.) Say it once more.
\end{tcolorbox}
\section[\texorpdfstring{\bn{ব্যবসা} (bæbsā)}{ব্যবসা (bæbsā)}]{\bn{ব্যবসা} (bæbsā) — business}

\label{lesson:BN-C207-bbsa}



\begin{quote}
Before the new one: say the Bengali for a shopkeeper, then the Bengali for a customer.
\end{quote}

\subsection*{You'll want to know: \bn{ব্যবসা}}
\textbf{\bn{ব্যবসা}} — \emph{bæbsā} — \textquotedblleft{}business\textquotedblright{}.

\begin{cousinweb}
\textbf{\bn{ব্যবসা}} (\emph{bæbsā}) is a tadbhava, a word inherited from Sanskrit and reshaped by sound change. It comes from \emph{vyavasāya}, \textquotedblleft{}resolve; trade\textquotedblright{} (CDIAL 11831). CDIAL lists Bengali \emph{byabsā}, \textquotedblleft{}business\textquotedblright{}, and notes the pull of the learned \emph{byabasāyī}.
\end{cousinweb}

\begin{grammarlens}[title={What you've built}]
One more word for this chapter.
\end{grammarlens}

\subsection*{Your turn}
\begin{itemize}
\raggedright
  \item \emph{Say it:} \emph{bæbsā}
  \item \emph{Say it:} \emph{bæbsā}, once more
  \item \emph{Say it:} say \emph{bæbsā}
  \item \emph{Recall:} say the Bengali for a shopkeeper, then the Bengali for a customer, then say \emph{bæbsā} again
\end{itemize}

\begin{tcolorbox}[breakable,colback=teal!4,colframe=teal!35!black,title={Before you move on}]
What does \bn{ব্যবসা} mean? (\textquotedblleft{}Business\textquotedblright{}.) Say it once more.
\end{tcolorbox}
\section[\texorpdfstring{\bn{কোম্পানি} (kompāni)}{কোম্পানি (kompāni)}]{\bn{কোম্পানি} (kompāni) — a company}

\label{lesson:BN-C207-kompani}



\begin{quote}
Before the new one: say the Bengali for a customer, then the Bengali for business.
\end{quote}

\subsection*{You'll want to know: \bn{কোম্পানি}}
\textbf{\bn{কোম্পানি}} — \emph{kompāni} — \textquotedblleft{}a company\textquotedblright{}.

\begin{cousinweb}
\textbf{\bn{কোম্পানি}} (\emph{kompāni}) is a loan from English \emph{company}.
\end{cousinweb}

\begin{grammarlens}[title={What you've built}]
One more word for this chapter.
\end{grammarlens}

\subsection*{Your turn}
\begin{itemize}
\raggedright
  \item \emph{Say it:} \emph{kompāni}
  \item \emph{Say it:} \emph{kompāni}, once more
  \item \emph{Say it:} say \emph{kompāni}
  \item \emph{Recall:} say the Bengali for a customer, then the Bengali for business, then say \emph{kompāni} again
\end{itemize}

\begin{tcolorbox}[breakable,colback=teal!4,colframe=teal!35!black,title={Before you move on}]
What does \bn{কোম্পানি} mean? (\textquotedblleft{}A company\textquotedblright{}.) Say it once more.
\end{tcolorbox}
